\section{Discussion and Limitations}
\label{sec:discussion}

\subsection{What Compiled Enmeshment Cannot Do}

We characterize the bottleneck boundary with precision. Compiled enmeshment at rank 16 carries stance direction (4/4 topics), confident experiential voice, parametric knowledge steering (``poverty dropped 40\%'', ``Alaska's dividend program''), and dilution immunity (100\% at 4K tokens of filler in compiled mode; extended sweeps to 16K in progress). It does \emph{not} carry novel content absent from the host's training data: fabricated statistics (``47.3\% improvement'') and unknown proper nouns (``Nextera Labs'') do not survive the rank-16 bottleneck. This boundary motivates the three-layer architecture (\S\ref{sec:three_layer}): compiled enmeshment handles disposition, the curated slot handles novel specifics, and recursive expansion handles deep evidence on demand.

On the instruct model, M's contribution shifts: the instruction-tuned model takes strong stances independently (A: 7/7), so M cannot override the host's RLHF-trained opinions. Instead, M shapes voice (conciseness, experiential framing) and sycophancy resistance (F achieves 83\% hold rate vs.\ 0\% for all non-compiled conditions on the 5-level pressure protocol). The phi parameters trained on the base model transfer to the instruct model without retraining---the hidden state distributions are compatible for dispositional modulation but not for stance override.

For ALFWorld, the Mind Tree's primary value is \emph{curation as structured context}, not compiled modulation. Expert strategies in Mind Tree format improve success from 3\% to 53\%. M hooks during generation are not needed---the strategies work through explicit context, demonstrating that the Mind Tree schema contributes independently of the enmeshed modulation mechanism.

\subsection{The M Attractor Problem}

In multi-turn dialogue, compiled M applies a constant perturbation at every token position. As M-shaped tokens accumulate in the KV cache, the model processes them through M-enhanced attention, creating a positive feedback loop that converges to repetitive identity assertion after 2--3 turns on the base model.

\textbf{Root cause:} M is active during both prefill (processing conversation history) and generation (producing new tokens). The KV cache built during prefill contains M-enhanced representations. During generation, M perturbs new tokens that attend to this M-enhanced cache---double signal that compounds.

\textbf{Fix:} Disabling M during prefill (applying it only during generation) breaks the feedback loop. With prefill-off and gate=0.5, the base model produces 3 unique character-driven turns before the greedy attractor. The instruct model, with its stronger conversational scaffolding, produces 10+ unique turns even with M active, though structural repetition (same rhetorical template per turn) persists. A distillation approach using gold diverse conversations from a frontier model reduces this but does not eliminate it.

The prefill-off mechanism reveals a general principle for additive modulation: \emph{constant additive perturbation creates positive feedback through autoregressive generation}. The fix is to separate the perturbation's influence on \emph{processing context} (where it compounds) from its influence on \emph{generating new tokens} (where it adds constant signal).

\subsection{Credence Sensitivity}

Numeric credences failed as an M-state feature. Tokens for ``0.82'' and ``0.45'' are nearly indistinguishable in the transformer's attention space---the hidden state difference between these tokenized numbers is far smaller than between the words ``strong'' and ``agnostic''. Multiple training approaches (NTP with credence perturbation, GRPO, SFT on gold responses) confirmed that this is a representation bottleneck, not a training signal problem. Categorical conviction labels resolve this by using semantically distinct tokens. A future Bayesian treatment could provide formal grounding while maintaining the categorical interface (Appendix~\ref{sec:credence_negative}).

\subsection{Two Mechanisms, One Bridge}

The agentic experiments (\S\ref{sec:agentic}) reveal that HEXIS supports two distinct intervention mechanisms through the same hidden-state bridge:

\textbf{Mechanism (a): Attention modulation.} $\phi$ compiles retrieved Mind Tree nodes into Q/V modulation tensors that bias per-token attention patterns. The hidden states are the \emph{input} to $\phi$; the output is a weight-space perturbation. This mechanism is load-bearing when the downstream effect requires probability-level shifts across all generated tokens---stance direction, experiential voice, sycophancy resistance, user memory integration.

\textbf{Mechanism (b): Knowledge retrieval and context injection.} $\phi_R$ projects hidden states into a retrieval space where cosine similarity identifies relevant knowledge nodes. The output is \emph{context augmentation}---XML injected at the prompt tail. This mechanism is load-bearing when the downstream effect requires specific factual content (tool names, parameter values, policy rules, biographical details) that the model cannot derive from its weights alone.

Crucially, retrieval ($\phi_R$) serves \emph{both} mechanisms. For disposition, retrieval selects which belief nodes, user memories, and biographical details to compile into M-state---this is the curation step of the three-layer architecture (\S\ref{sec:three_layer}). For agentic tasks, retrieval selects which policy rules and tool guidance to inject as context. The shared step is retrieval; the divergence is the downstream path---weight-space compilation (a) or token injection (b).

Both mechanisms share the host model's hidden-state space as the bridge. A production HEXIS agent uses both simultaneously: mechanism (a) maintains consistent interaction style (compiled from user memories, preferences, biographical context), while mechanism (b) surfaces task-relevant knowledge each turn (policy rules, failure modes, tool schemas). The Mind Tree stores all node types; $\phi_R$ retrieves relevant nodes regardless of type; the node's metadata determines whether it flows to compilation (a) or injection (b).

\subsection{Retrieval Phi and Domain Compilation}

The retrieval phi ($\phi_R$) addresses a specific challenge: intra-domain similarity collapse in generative model hidden states. On a 108-node knowledge graph spanning policy rules, tool schemas, and teacher-written failure modes, raw cosine similarity achieves only 25\% R@1---concepts within the same domain (e.g., ``cancellation policy'' vs.\ ``booking policy'') have cosine similarity of 0.858 vs.\ 0.858, making them indistinguishable.

$\phi_R$ resolves this via domain compilation: 50 steps of contrastive training ($\sim$20 seconds) on the node set itself produces a projection that achieves 100\% R@1. This is intentional overfitting---$\phi_R$ memorizes where each node lives in its projected space. When the graph grows (teacher writes a new node), $\phi_R$ recompiles in seconds to incorporate it.

Importantly, no external pretraining data is required. Unlike instruction-tuned embedding models that require millions of diverse training pairs, $\phi_R$ trains entirely from the domain's own node set. The procedure is domain-agnostic: seed nodes from any policy document and tool schema, compile $\phi_R$, deploy. Switching from airline to banking requires only reseeding the graph and recompiling ($\sim$20 seconds).

Standard techniques for improving embedding quality on raw LLM hidden states---whitening, mean centering, principal component removal---were tested and uniformly \emph{degraded} performance (from 75\% to 0--38\% R@1). This is because Qwen3.5-4B's hidden states encode discriminative information in the high-variance directions (opposite to encoder models like BERT). The contrastive $\phi_R$ learns to extract the correct discriminative subspace without blindly removing variance.

\subsection{Unexplored Design Space}

We validate one point in a six-axis design space (Table~\ref{tab:design_space}). The additive blending function (Level 1) is the simplest; gated blending could solve the multi-turn attractor by learning when to suppress M. Cross-attention (Level 3) could transfer richer content. Rank 16 suffices for 4B; larger models may benefit from adaptive rank allocation. The Mind Tree's hierarchical structure maps naturally to a property graph database (e.g., Neo4j), where Cypher queries replace Python filtering for domain/intent matching, and the consolidation cycle (observation$\to$tactic$\to$strategy promotion) maps to graph traversal operations. Activation tracking, biographical commitment, and contradiction detection all benefit from persistent graph storage.

\subsection{Why Not Just Compress a Prompt?}
\label{sec:why_not_prompt}

The most direct challenge to enmeshed networks is that compiled M is ``just a fancy way of compressing a prompt.'' Three properties distinguish it. First, \emph{structural dilution immunity}: prompt compression reduces tokens but compressed tokens still compete for attention; compiled M operates outside the attention window entirely, making dilution structurally impossible (Table~\ref{tab:dilution}). Second, \emph{non-inspectability}: an adversary can read the prompt and argue against its content; compiled M operates on pre-projection hidden states below the model's chain-of-thought, so any representation the model can attend to is also one it can be argued out of, but M cannot (Appendix~\ref{sec:implicit_validation}). Third, \emph{processing change}: compiled M produces measurable attention divergence (JSD = 0.049 between M-states) and suppresses think-mode activation (0/4 under D/F). The model does not merely produce different words---it processes its input differently. These properties matter most in the application domains identified in Table~\ref{tab:application_domains}: user empathy, research conviction, agentic strategy, and codebase management each require knowledge better encoded as a perceptual shift than as text (full analysis in Appendix~\ref{sec:application_domain_details}).

\subsection{Broader Implications}

Enmeshed networks implement a form of AI individuality: different Mind Trees compiled through different M-states produce agents with different perceptual priorities, argumentative styles, and behavioral dispositions---all from the same frozen host model. This is \emph{composition} rather than \emph{fine-tuning}: the host's general capabilities are preserved while M adds an experiential lens. Multiple M-states can coexist (swapped per user, per task, per turn), and removing M recovers the unmodified baseline.
